%%
%% ACM Journal Article — Alagang MMC
%% Author: Augustus Clark Raphael P. Rodriguez
%%
%% SUBMISSION MODE: Use \documentclass[manuscript, screen, review]{acmart}
%% CAMERA-READY:   Use \documentclass[acmsmall]{acmart}
%%
\documentclass[manuscript, screen, review]{acmart}

%% ============================================================
%% PREAMBLE: Rights & Journal Metadata
%% Fill these in once you receive the rights form from ACM.
%% Leave as-is during drafting/review.
%% ============================================================
\setcopyright{acmcopyright}
\copyrightyear{2026}
\acmYear{2026}
\acmDOI{XXXXXXX.XXXXXXX}            % Provided by ACM after acceptance

%% Which journal are you submitting to? Replace JACM with the correct
%% journal code. Common examples: TOCHI, TOSEM, HEALTH, PACMHCI.
\acmJournal{JACM}
\acmVolume{0}                        % Filled in by ACM
\acmNumber{0}                        % Filled in by ACM
\acmArticle{0}                       % Filled in by ACM
\acmMonth{0}                         % Filled in by ACM

%% ============================================================
%% DOCUMENT BODY
%% ============================================================
\begin{document}

%% ------------------------------------------------------------
%% TITLE
%% The optional argument is the short title shown in page headers.
%% Keep the short title under ~50 characters.
%% ------------------------------------------------------------
\title[CeHRes Roadmap for Care Coordination in the Philippines]{Designing for Digital Health: 
A CeHRes Roadmap Application for Care Coordination in the Philippines}

%% ------------------------------------------------------------
%% AUTHORS & AFFILIATIONS
%% Each \author block must have a matching \affiliation block.
%% \email and \orcid are optional but encouraged.
%% For the journal version, do NOT use \authornote for "equal contribution"
%% unless it genuinely applies — it changes the footnote marker behavior.
%% ------------------------------------------------------------
\author{Augustus Clark Raphael P. Rodriguez}
\email{aprodriguez7@up.edu.ph}
% \orcid{XXXX-XXXX-XXXX-XXXX}       % Add your ORCID if you have one
\affiliation{%
  \institution{University of the Philippines Manila}
  \department{Department of Physical Sciences and Mathematics}
  \city{Manila}
  \country{Philippines}
}

\author{Marbert John C. Marasigan}
\email{mcmarasigan2@up.edu.ph}    
\affiliation{%
  \institution{University of the Philippines Manila}
  \department{Department of Physical Sciences and Mathematics}
  \city{Manila}
  \country{Philippines}
}

%% Short author list for page headers (use "et al." if >2 authors)
\renewcommand{\shortauthors}{Rodriguez and Marasigan}

%% ------------------------------------------------------------
%% ABSTRACT
%% Target: 150–250 words. Should cover: problem, method, key results,
%% and significance. This is NOT your thesis abstract — condense it.
%% A journal abstract must be self-contained (no citations).
%% ------------------------------------------------------------
\begin{abstract}
  % TODO: Write a condensed, self-contained abstract covering:
  %   1. The problem: process hurdles in MMC's patient care coordination
  %   2. The gap: limited integration and disjoint systems
  %   3. Your approach: Human-Centered Design via the CeHRes Roadmap
  %   4. What you built: Alagang MMC (describe system briefly)
  %   5. Key results: UEQ scores, usability findings, stakeholder outcomes
  %   6. Significance: contribution to digital health in Philippine context
  %
  % Example structure (replace with real content):
  % "Patient care coordination in large Philippine hospitals faces [problem].
  %  Existing systems [gap]. We designed and developed Alagang MMC, a
  %  [description], using the CeHRes Roadmap to ensure stakeholder alignment
  %  throughout all phases. Evaluation via the User Experience Questionnaire
  %  revealed [results]. These findings demonstrate [significance]."
\end{abstract}

%% ------------------------------------------------------------
%% CCS CONCEPTS (MANDATORY for papers over 2 pages)
%% Generate the XML block at: https://dl.acm.org/ccs
%% Search for concepts like "Human-centered computing", "Health informatics"
%% Paste the generated block here and add the matching \ccsdesc lines below.
%% ------------------------------------------------------------
\begin{CCSXML}
% TODO: Replace this entire block with the output from https://dl.acm.org/ccs
% Suggested concepts to search for:
%   - Human-centered computing > User interface design
%   - Applied computing > Health informatics
%   - Human-centered computing > Participatory design
%   - Applied computing > Health care information systems
<ccs2012>
 <concept>
  <concept_id>REPLACE_WITH_GENERATED_ID</concept_id>
  <concept_desc>Human-centered computing~User interface design</concept_desc>
  <concept_significance>500</concept_significance>
 </concept>
 <concept>
  <concept_id>REPLACE_WITH_GENERATED_ID</concept_id>
  <concept_desc>Applied computing~Health informatics</concept_desc>
  <concept_significance>300</concept_significance>
 </concept>
</ccs2012>
\end{CCSXML}

% TODO: Match these \ccsdesc lines to your CCSXML block above
\ccsdesc[500]{Human-centered computing~User interface design}
\ccsdesc[300]{Applied computing~Health informatics}

%% ------------------------------------------------------------
%% KEYWORDS (MANDATORY for papers over 2 pages)
%% 4–6 keywords, comma-separated, lowercase except proper nouns.
%% ------------------------------------------------------------
\keywords{human-centered design, digital health, care coordination,
  CeHRes Roadmap, health information systems, Philippines}

%% ------------------------------------------------------------
%% RECEIVED DATES — filled in by journal editors, leave as placeholders
%% ------------------------------------------------------------
\received{20 January 2026}
\received[revised]{MONTH DAY, YEAR}
\received[accepted]{MONTH DAY, YEAR}

%% This command generates the title block, author list, abstract, etc.
\maketitle

%% ============================================================
%% SECTION 1: INTRODUCTION
%% Journal intro is tighter than a thesis intro. Target: ~600–900 words.
%% Structure: context → problem → gap in literature → your contribution
%% → paper organization (1 short paragraph at the end of the intro).
%% ============================================================
\section{Introduction}
% TODO: Adapt from your thesis Background of the Study + Statement of the Problem.
% Compress into a single narrative. Do NOT list subsections here — journal
% introductions are typically a flowing discussion, not subdivided.
%
% Cover:
%   - Digital health landscape in the Philippines and MMC's context
%   - The three process hurdles (limited accessibility, lack of integration,
%     disjoint systems) — weave them into the narrative rather than bullets
%   - HCD and CeHRes Roadmap as your response to the gap
%   - Your three research questions (can be stated as numbered RQs inline)
%   - A closing "roadmap" paragraph: "The rest of this paper is organized
%     as follows: Section 2 reviews related work... Section 3 describes..."


%% ============================================================
%% SECTION 2: RELATED WORK / BACKGROUND
%% Situate your work in prior literature. ~800–1200 words.
%% Organize by theme, not chronologically.
%% ============================================================
\section{Related Work}
% TODO: Synthesize literature from your thesis into thematic subsections.
% Suggested subsections:

\subsection{Digital Health Systems in Developing Countries}
% Context on HIS in the Philippines (Sucaldito et al., Guzman, etc.)
% What works, what doesn't, what gaps remain.

\subsection{Human-Centered Design in Healthcare}
% HCD philosophy and its healthcare applications (Leary, Landeiro et al.)
% CeHRes Roadmap specifically — describe its phases and why it fits.

\subsection{Patient Care Coordination Systems}
% Existing coordination systems, their limitations.
% How Alagang MMC addresses the gap.


%% ============================================================
%% SECTION 3: METHODOLOGY
%% Explain what you did and how. ~1000–1500 words.
%% Must be reproducible: a reader should understand how to replicate your study.
%% ============================================================
\section{Methodology}

\subsection{Research Design}
% State your overall approach: action research / design science research /
% user-centered iterative development — whichever framing fits your work.
% Briefly justify why this design is appropriate.

\subsection{CeHRes Roadmap Application}
% Walk through each phase of the CeHRes Roadmap as you applied it:
%   - Contextual Inquiry (stakeholder interviews: Dr. Rodriguez, Dr. Javier)
%   - Value Specification
%   - Design
%   - Operationalization
%   - Summative Evaluation
% Include your DFD, UCD, and activity diagrams where appropriate
% (use \includegraphics with images from sp-manuscript/images/).

\subsection{System Architecture}
% Describe the technical design of Alagang MMC:
%   - Key modules: Clinician Profile Manager, Appointment Management,
%     Clinician Directory, Account Creation, Appointment Booking
%   - Technology stack
%   - Include your ERD and pipeline diagram if relevant

\subsection{Evaluation Approach}
% Describe how you evaluated: UEQ, sprint-based usability testing,
% participants (patients, secretaries, clinicians), protocol.


%% ============================================================
%% SECTION 4: RESULTS
%% Present findings objectively. ~800–1200 words.
%% Use tables and figures to present data efficiently.
%% ============================================================
\section{Results}

\subsection{Usability and User Experience Findings}
% Present UEQ results (benchmark comparison, scale scores).
% Reference your benchmark figure and scale structure figure.
% State which aspects scored above/below benchmark and what that means.

\subsection{Stakeholder Feedback}
% Qualitative findings from sprint reviews, bug reports,
% and clinician/patient interviews.
% Tie back to your three research questions.

\subsection{System Outcomes}
% Any quantitative metrics on system usage or workflow improvement.


%% ============================================================
%% SECTION 5: DISCUSSION
%% Interpret your results, link to literature, acknowledge limitations.
%% ~600–900 words.
%% ============================================================
\section{Discussion}

\subsection{Implications for Digital Health in the Philippines}
% What do your results mean for HIS adoption in developing-country hospitals?
% Connect to Sucaldito et al., Livieri et al., Libadia et al.

\subsection{HCD and the CeHRes Roadmap as a Framework}
% How did stakeholder-driven development influence outcomes (RQ3)?
% Compare to Landeiro et al. or similar HCD-in-healthcare work.

\subsection{Limitations}
% Be honest: single-site study, sample size, generalizability,
% any technical constraints. This strengthens credibility.


%% ============================================================
%% SECTION 6: CONCLUSION
%% ~200–350 words. No new information. Summarize contribution,
%% answer your three RQs, and state future work.
%% ============================================================
\section{Conclusion}
% TODO: Summarize:
%   - What you built and why
%   - Direct answers to your three research questions
%   - Significance of using HCD/CeHRes for care coordination
%   - Future work: expansion beyond MMC, integration with iHIMS,
%     longitudinal adoption studies, etc.


%% ============================================================
%% ACKNOWLEDGMENTS
%% Thank anyone who helped but is not an author.
%% Do NOT include your adviser here if they are a co-author above.
%% ============================================================
\begin{acks}
  % TODO: Acknowledge funding sources, research participants (can be kept
  % anonymous as "the clinicians and staff of MMC who participated"),
  % and any institutional support.
  %
  % Example: "The authors thank the clinicians, secretaries, and patients
  % of Makati Medical Center who participated in this study. This work
  % was supported by [funding source, if any]."
\end{acks}

%% ============================================================
%% BIBLIOGRAPHY
%% The .bib file below must be in the same directory as this .tex file.
%% The ACM-Reference-Format.bst style is provided by acmart.
%% ============================================================
\bibliographystyle{ACM-Reference-Format}
\bibliography{references}

\end{document}
